\usepackage[ngerman]{babel}					% TeX-Ausgabe-Sprache
\usepackage[backend=biber, style=alphabetic]{biblatex}	% Paket fuer Literaturverzeichnisse
\usepackage{blindtext}						% Blindtext um den Grauwert der Schrift zu pruefen
\usepackage{booktabs}						% Linien für Tabellenpaket
\usepackage{csquotes}						% Kontextsensitive Zitate (Wird von biblatex gefordert)
\usepackage{graphicx}						% Einbinden von Grafiken
\usepackage{listings} 						% Quellcode darstellen
\usepackage{setspace}						% Durchschuss
\usepackage{siunitx}						% Schoene Einheiten
\usepackage{subcaption}						% Aufgeteilte Bilder
\usepackage{xcolor}							% Farben
\usepackage{textpos}						% Zur absoluten Positionierung von Textblöcken (Titelseite)
\usepackage{geometry}						% Zur freien Anpassung der Dokumenteneigenschaften (Titelseite)

% DAS HYPERREF- UND GLOSSARIES-PAKET SOLLTE ALS LETZTE EINGEBUNDEN WERDEN
\usepackage{hyperref}						% Hyperlinks im Dokument
\usepackage{xurl}							% Erstellt Zeilenumbrüche bei URLs
\usepackage[toc, acronyms, automake]{glossaries}	% Glossar erstellen

% EINSTELLUNGEN

\hypersetup{
			unicode		=false,				% Ermöglicht nicht-latin Schriftzeichen
			pdftitle	={\theTitle},		% Titel der Arbeit für die Dokumenteninfo
			pdfauthor	={\theAuthor},		% Name des Autors
			pdfsubject	={\theAbstractEng},	% Abstract der Arbeit
			pdfcreator	={\theAuthor},		% Name des Erstellers
			pdfproducer	={\theAuthor},		% Name des Produzenten
			pdfkeywords	={\theKeywordsEng},	% Schlüsselwörter des Dokuments
			pdfdisplaydoctitle = true,		% Zeigt den Titel in der Titelzeile statt Dateinamen
			pdfmenubar	=true,				% Menuleiste im PDF-Viewer anzeigen
			pdftoolbar	=true,				% PDF-Toolbar im PDF-Viewer anzeigen
			colorlinks	=true,				% farbiger Rahmen um die Links Schrift farbig
			linkcolor	=red,				% Farbe für die Links
			linktoc		=all,				% none,section,page,all
			citecolor	=green,				% Farbe für die Quelllinks
			filecolor	=magenta,			% Farbe für Dateilinks
			urlcolor	=cyan,				% Farbe für URL-Links (Web, Mail)
%			hidelinks,						% Hyperlinks unmarkiert
			%	linkbordercolor	={1 0 0}, 	% Farbe für den Rahmen um die Links
			%	citebordercolor	={0 1 0}, 	% Farbe für den Rahmen um die Links für Zitate
			%	urlbordercolor	={0 1 1}, 	% Farbe für den Rahmen um die URL-Links
}

\sisetup{
			locale = DE ,
			per-mode = symbol
}

\addtokomafont{descriptionlabel}{\rmfamily}		% glossaries-Verzeichnisse mit Serifen